\clearpage

\lehead[]{\normalfont\sffamily\hspace*{-2.00cm}\textcolor{white}{\colorbox{lightblue}{\parbox[c][0.70cm][b]{1.60cm}{
\makebox[1.60cm][r]{\thechapter}\\ \makebox[1.60cm][r]{ÜBUNG}}}}\hspace{0.17cm}\textcolor{lightblue}{\chaptertitle}}
\rohead[]{\textcolor{lightblue}{\chaptertitle}\normalfont\sffamily\hspace*{0.17cm}\textcolor{white}{\colorbox{lightblue}{\parbox[c][0.70cm][b]{1.60cm}{\thechapter\\
ÜBUNG}}}\hspace{-2.00cm}}
%\chead[]{}
\rehead[]{\textcolor{lightblue}{AvHG, Inf, My}}
\lohead[]{\textcolor{lightblue}{AvHG, Inf, My}}

\section{Erste Übung -- erster Teil}

Für die folgende Übung braucht ihr nur Blatt und Stift. Wer möchte, darf
zusätzlich einen einfachen Taschenrechner benutzen (die Grundrechenarten
$+$, $-$, $\cdot$, $:$ genügen).

\paragraph{Das Problem:}
Die Quadratwurzel einer beliebigen positiven Zahl soll auf zwei Nachkommastellen
genau bestimmt werden, \emph{ohne}(!) die Wurzelfunktion des Taschenrechners zu
benutzen.

\paragraph{Euer Auftrag:}
\begin{compactenum}[1.]
\item Setzt euch in Zweier- oder Dreiergruppen zusammen.
\item Überlegt euch anhand eines konkreten Beispiels (z.\,B.\ $\sqrt{20}$), wie
ihr systematisch vorgehen könnt, um dem exakten Ergebnis immer näher zu kommen.
\item Formuliert aus eurem Vorgehen eine allgemeine, eindeutige
\textbf{Schritt-für-Schritt-Anleitung} (einen \emph{Algorithmus}), die für
jede positive Zahl funktioniert.
\end{compactenum}

\medskip
\noindent Dabei spielt es zunächst noch keine Rolle, wie kurz oder elegant euer
Lösungsweg ist -- Hauptsache, er funktioniert zuverlässig!

\medskip
\noindent \textbf{Zeit:} 15 Minuten \\
\textbf{Sicherung:} Schreibt eure Anleitung so übersichtlich auf, dass ihr sie
anschließend der Klasse vorstellen und an einem neuen Beispiel testen könnt.
